\section{Conventions and the nonnegative scalar-curvature reduction}
\label{sec:reductions}

We record the degree conventions for the comparison construction and reduce
nonnegative scalar curvature to the positive and Ricci-flat cases used in the
corollaries. All manifolds and metrics are smooth. A closed manifold is compact
and has no boundary. Unless stated otherwise, manifolds are oriented.
Fundamental classes are taken with the indicated coefficients; we use the same
symbol $[M]$ for
an integral fundamental class and its images under extension of coefficients.
The Laplacian is $\Delta=\diver\nabla$, and a round sphere has positive scalar
curvature. For a hypersurface with unit normal $\nu$, our second fundamental
form is $\mathrm{II}(X,Y)=\langle\nabla_X\nu,Y\rangle$ and its mean curvature
is its trace. These conventions fix the signs in the graph calculations below.
We give $S^d$ its unit round metric and write
$\T^b=(\R/2\pi\Z)^b$ for the standard flat torus.

\subsection{Fundamental classes and degree}

Throughout the proof, $\Gamma=\pi_1(M)$ and $c_M:M\to B\Gamma$
have the meaning fixed in \cref{sec:introduction}. No finite-dimensional
model for $B\Gamma$ is assumed. Disconnected closed manifolds are treated
one component at a time; compactness leaves only finitely many components.

For a possibly noncompact oriented $d$-manifold $Y$, a smooth map
$F:Y\to S^d$ that is equal to one value $*\in S^d$ outside a compact set has
a well-defined degree. One may compute it as the signed number of preimages
of any regular value different from $*$. Those preimages are contained in a
compact set. Equivalently, pull back a compactly supported top-degree class on
$S^d\setminus\{*\}$ and evaluate it on the locally finite fundamental class
of $Y$. This definition also permits disconnected $Y$: only finitely many
components can meet the compact region on which $F$ is nonconstant. We use
\emph{constant outside a compact set} in this sense, rather than assigning a
linear meaning to compact support for a sphere-valued map.

We also use the ordinary oriented local degree $\deg(G,\Omega,0)$ for a
smooth map from an oriented $q$-manifold to $\R^q$, where $\Omega$ is a
precompact open set and $0\notin G(\partial\Omega)$. Smoothness of
$\partial\Omega$ is unnecessary. A homotopy whose boundary values avoid zero
preserves this degree. More generally, if they avoid a fixed ball, composition
with a degree-one map $\R^q\to S^q$ constant outside that ball gives a
sphere-valued map of the same degree, extending constantly across the boundary
of $\Omega$.

\begin{lemma}[Simplicial volume detects rational essentiality]
\label{lem:essentiality}
If a closed connected oriented manifold $M^n$ has positive real simplicial
volume, then $(c_M)_*[M]\ne0$ in $H_n(B\Gamma;\Q)$.
\end{lemma}

\begin{proof}
The bounded-cohomology mapping theorem identifies real bounded cohomology of
$M$ with that of $B\Gamma$ by pullback along $c_M$. The duality between
bounded cohomology and the $\ell^1$-seminorm implies that
$\|M\|>0$ supplies a bounded class whose ordinary cohomology image evaluates
nontrivially on $[M]$. These are the mapping and duality results of
\cite{Gromov1982}; neither requires a spin structure.

If $(c_M)_*[M]$ vanished rationally, its real image would vanish. Every bounded
class pulled back from $B\Gamma$ would then evaluate trivially on $[M]$, a
contradiction. In dimension one, real bounded cohomology already vanishes and
the assertion is vacuous; the direct circle argument below also suffices.
\end{proof}

\subsection{The scalar-flat alternative}

\begin{lemma}[Positive scalar curvature or Ricci flatness]
\label{lem:scalar-reduction}
Let $M$ be closed and connected, not necessarily oriented, and suppose that
$g$ has nonnegative scalar curvature. Either $M$ admits a metric of strictly positive scalar curvature,
or $g$ is Ricci-flat. In the Ricci-flat case, if $M$ is oriented and has
positive dimension, then $\|M\|=0$.
\end{lemma}

\begin{proof}
Run the short-time Ricci flow starting at $g$. Short-time existence on closed
manifolds and the scalar evolution equation give
\begin{equation}\label{eq:scalar-evolution}
 \partial_t\Scal=\Delta\Scal+2|\Ric|^2.
\end{equation}
See \cite[Theorem~5.2.1 and Proposition~2.5.4]{Topping2006} for these
statements in arbitrary dimension, and \cite{Hamilton1982} for the
Ricci-flow framework. The maximum principle
preserves nonnegativity. If the initial scalar curvature is positive at one
point, the strong maximum principle makes it positive everywhere at any
sufficiently small later time. If it is initially identically zero but the
Ricci tensor is nonzero at a point, \eqref{eq:scalar-evolution} gives a
strictly positive initial time derivative there. At a sufficiently small positive time $t_1$, scalar curvature is
therefore positive at that point. Applying the strong maximum principle
from $t_1$ produces a metric of positive scalar curvature everywhere at
a later time. Thus failure of this alternative
forces $\Ric_g=0$.

For the final assertion, suppose that $g$ is Ricci-flat, that $M$ is
oriented, and that $n=\dim M>0$. We use the orbit-packing argument
of Milnor \cite[Theorem~1]{Milnor1968}. Its complete universal Riemannian cover
$(\widetilde M,\widetilde g)$ has nonnegative Ricci curvature. Fix
$x\in\widetilde M$ and a finite generating set of $\Gamma$. There is a
constant $D$ such that
\[
 d(x,\gamma x)\le D|\gamma|
 \qquad(\gamma\in\Gamma),
\]
where $|\gamma|$ is word length. Choose a positive radius $r$ for which the
balls $B(\gamma x,r)$, $\gamma\in\Gamma$, are disjoint. To choose $r$, take a sufficiently small metric ball about the image of
$x$ contained in an evenly covered neighborhood. The radius-$r$ balls
about its lifts lie in distinct sheets and are therefore disjoint. All these
balls have the same positive volume, say $v_r$. The orbit balls with
$|\gamma|\le k$ lie in $B(x,Dk+r)$. Bishop--Gromov comparison \cite[Theorem~4.1]{Wei2007} therefore gives
\[
 \#\{\gamma:|\gamma|\le k\}\,v_r
 \le \Vol B(x,Dk+r)\le \omega_n(Dk+r)^n.
\]
Thus $\Gamma$ has polynomial word growth. Write
$b(k)=\#B_\Gamma(1,k)$. There are integers $k_j\to\infty$ such that
$b(k_j+1)/b(k_j)\to1$. Otherwise, some $\varepsilon>0$ would satisfy
$b(k+1)/b(k)\ge1+\varepsilon$ for all sufficiently large $k$, forcing
exponential growth. For each generator $s$,
\[
 \#\bigl(sB_\Gamma(1,k_j)\mathbin\triangle B_\Gamma(1,k_j)\bigr)
 \le2\bigl(b(k_j+1)-b(k_j)\bigr).
\]
These balls form a F\o lner sequence, so $\Gamma$ is amenable. This
uses only the growth bound, without a structural classification of groups
of polynomial growth. Real bounded
cohomology of an amenable group vanishes in positive degrees
\cite[Sections~3.0--3.1]{Gromov1982}. The mapping
theorem and duality used in \cref{lem:essentiality} now imply
$\|M\|=0$.
\end{proof}

Every closed connected oriented one-manifold is a circle. Its self-cover of
degree two and functoriality of the simplicial seminorm give
$2\|S^1\|\le\|S^1\|$, hence $\|S^1\|=0$. Consequently, to prove the
simplicial-volume theorem in every positive dimension, it remains to show
that a positive-scalar-curvature manifold of dimension at least two is
rationally inessential. Compactness lets us write its scalar lower bound as
$\Scal_M\ge\kappa>0$.
